\chapter{Barrett's Esophagus}

\section*{Definition and Overview}
Barrett's esophagus is a condition in which the flat pink lining of the lower esophagus (the swallowing tube connecting the throat to the stomach) is replaced by a red, velvety lining that resembles the lining of the stomach or intestine. The change is a response to long-term damage from acid reflux. The importance of Barrett's esophagus is that it slightly increases the risk of developing a type of esophageal cancer called adenocarcinoma, so it requires monitoring even though most people with the condition never develop cancer.

The condition is usually diagnosed on endoscopy and confirmed by biopsy. It does not cause symptoms of its own; instead, people are generally aware of the reflux symptoms that led to the change. Because the cancer risk, although small, is real, regular surveillance is recommended for most people with the diagnosis.

\section*{Epidemiology}
Barrett's esophagus is found in roughly 1 in 20 to 1 in 50 adults, and it is more common in men than women. It is most often diagnosed between the ages of 50 and 70 years. The condition is strongly linked to chronic acid reflux (gastro-oesophageal reflux disease, GORD); around 1 in 10 people with long-standing reflux develop Barrett's esophagus. People of European descent are affected more often than people of Asian or African descent, and rates have risen in recent decades, partly because of increased endoscopy use. Risk is higher in people with a family history of the condition.

\section*{Aetiology and Pathophysiology}
\begin{itemize}
    \item \textbf{Chronic acid reflux:} repeated exposure of the lower esophagus to stomach acid and bile damages the normal squamous lining
    \item \textbf{Metaplasia:} in response to injury, the damaged lining is replaced by columnar cells similar to those in the gut -- this change is the defining feature of Barrett's esophagus
    \item \textbf{Risk factors:} long-standing reflux, being male, older age, obesity (particularly around the abdomen), smoking, and a family history of Barrett's esophagus or esophageal cancer
    \item \textbf{Progression:} in a small minority, the metaplastic cells develop increasing abnormality (dysplasia) that can progress to adenocarcinoma over many years
\end{itemize}

\section*{Clinical Features}
Barrett's esophagus itself usually causes no symptoms; it is most often found incidentally during endoscopy performed for reflux symptoms. People with the condition commonly report:

\begin{itemize}
    \item Heartburn (a burning sensation behind the breastbone)
    \item Regurgitation of stomach contents into the throat or mouth
    \item Difficulty swallowing or a feeling of food sticking
    \item Chest discomfort or pain
    \item Chronic cough, hoarseness, or sore throat from reflux
\end{itemize}

\section*{Differential Diagnosis}
\begin{itemize}
    \item \textbf{Reflux oesophagitis:} inflammation of the esophagus from acid without Barrett's changes
    \item \textbf{Eosinophilic oesophagitis:} an allergic inflammatory condition of the esophagus
    \item \textbf{Oesophageal stricture or ring:} narrowing causing difficulty swallowing
    \item \textbf{Achalasia:} a nerve condition causing the esophagus to fail to empty
    \item \textbf{Oesophageal cancer} (adenocarcinoma or squamous cell carcinoma) -- must be excluded by biopsy
    \item \textbf{Oesophageal infection} (for example, candida or herpes) in immunocompromised people
\end{itemize}

\section*{When to Seek Medical Care}
See a doctor for persistent heartburn or reflux that is severe, frequent, or not responding to over-the-counter treatment, or for difficulty swallowing. Once Barrett's esophagus is diagnosed, attend regular surveillance as recommended.

\textbf{Call 000 (or local emergency services) for:}
\begin{itemize}
    \item Difficulty or pain on swallowing that is worsening
    \item Unintentional weight loss
    \item Vomiting blood or passing black, tarry stools
    \item Persistent chest pain, especially if it differs from your usual reflux pain
\end{itemize}

\section*{Diagnosis and Investigations}
\begin{itemize}
    \item \textbf{Endoscopy (gastroscopy):} a camera on a flexible tube is passed into the esophagus to inspect and photograph the lining
    \item \textbf{Biopsy:} small tissue samples are taken from the abnormal area and examined under a microscope to confirm Barrett's changes and detect any dysplasia
    \item \textbf{Surveillance endoscopy:} regular repeat gastroscopy every 2--5 years (or more often if dysplasia is found) to look for early cancer changes
    \item \textbf{Chromoendoscopy or high-definition endoscopy:} techniques that improve visualisation of suspicious areas during surveillance
    \item \textbf{pH monitoring (in selected cases):} a test to measure acid exposure in the esophagus when the diagnosis is unclear
\end{itemize}

\section*{Management}
\begin{itemize}
    \item \textbf{Acid suppression:} proton pump inhibitors (for example, omeprazole and pantoprazole) reduce acid production and help protect the esophagus
    \item \textbf{Lifestyle measures:} weight loss, avoiding late meals, raising the head of the bed, and avoiding alcohol, tobacco, and large fatty meals
    \item \textbf{Surveillance:} regular endoscopy to monitor the lining and detect dysplasia or early cancer
    \item \textbf{Endoscopic treatment:} for high-grade dysplasia or early cancer, techniques such as radiofrequency ablation or endoscopic mucosal resection can remove the abnormal tissue
    \item \textbf{Surgery:} anti-reflux surgery (fundoplication) may be considered in selected people, though it does not reliably eliminate the cancer risk
    \item \textbf{Follow-up:} long-term review by a gastroenterologist for everyone with the diagnosis
\end{itemize}

\section*{Complications}
\begin{itemize}
    \item Dysplasia -- precancerous changes in the Barrett's lining, graded as low or high grade
    \item Adenocarcinoma of the esophagus -- the risk is roughly 1 in 300 to 1 in 1000 per year, and it is higher with high-grade dysplasia
    \item Oesophageal stricture -- narrowing of the esophagus from scarring, causing difficulty swallowing
    \item Complications of endoscopy itself, such as bleeding or perforation (rare)
    \item Side effects of long-term acid-suppressing medicines, including altered nutrient absorption
\end{itemize}

\section*{Prognosis and Outcomes}
Most people with Barrett's esophagus never develop cancer, and the condition is often stable for many years. Low-grade dysplasia may regress, stay stable, or progress over time, while high-grade dysplasia carries a higher risk of progression. With surveillance and modern endoscopic treatment, early cancers can usually be cured. Life expectancy for most people with Barrett's esophagus is close to normal, and the main benefits of monitoring are early detection and treatment of precancerous or early cancerous changes.

\section*{Prevention and Self-Care}
\begin{itemize}
    \item Manage reflux well: take acid-suppressing medicines as prescribed and review them regularly
    \item Maintain a healthy body weight, since abdominal obesity increases reflux and Barrett's risk
    \item Eat smaller meals and avoid eating within 2--3 hours of bedtime
    \item Avoid alcohol, tobacco, and foods that trigger your reflux (spicy, fatty, or acidic foods)
    \item Raise the head of your bed slightly to reduce night-time reflux
    \item Attend all scheduled surveillance endoscopies
    \item Report new or worsening symptoms, such as difficulty swallowing or weight loss, to your doctor promptly
\end{itemize}

\section*{Sources and Further Reading}
The following sources provide reliable, up-to-date information on Barrett's esophagus.
\begin{itemize}
    \item NHS. \textit{Barrett's oesophagus.} https://www.nhs.uk/conditions/
    \item Mayo Clinic. \textit{Barrett's esophagus.} https://www.mayoclinic.org
    \item MSD Manual. \textit{Barrett Esophagus.} https://www.msdmanuals.com
    \item NICE. \textit{Barrett's oesophagus: surveillance.} https://www.nice.org.uk
    \item MedlinePlus. \textit{Barrett Esophagus.} https://medlineplus.gov
\end{itemize}
